% 2008 年江西卷理科第 16 题：正四棱柱容器内底部镶嵌同底的正四棱锥（顶点 P），水面恰好过 P；
%   图 1 为水平放置，图 2 为倒置后水面同样过 P。
% 来源 PDF：gaokaomath-src/普通高考/2008/2008江西理.pdf
% SHA-256：8d079d71fadf20911ed2fd71acfe9562b94eb962cfe5228fe7a376f3b599b239
% 原页：2008江西理.pdf 第 1 页（A3 三栏）第 3 栏，书页码 459。
% 绘制依据：棱柱高取 5、棱锥高取 3（正好一半略多，仅作示意）。斜投影下底面为 A(0,0)、B(3,0)、
%   C(4.2,0.9)、D(1.2,0.9)。图 1：棱锥底面贴底、顶点 P 在底面中心上方高 3 处，水面过 P；
%   图 2：容器倒置，棱锥顶点朝下，水面同样过 P。水用浅灰填充。
\begin{tikzpicture}[scale=0.52, semithick, every node/.style={inner sep=1.2pt, font=\small}]

  % ---------- 图 1：正放 ----------
  \begin{scope}[shift={(0,0)}]
    \coordinate (A)  at (0,0);
    \coordinate (B)  at (3,0);
    \coordinate (C)  at (4.2,0.9);
    \coordinate (D)  at (1.2,0.9);
    \coordinate (A1) at (0,5);
    \coordinate (B1) at (3,5);
    \coordinate (C1) at (4.2,5.9);
    \coordinate (D1) at (1.2,5.9);
    \coordinate (P)  at (2.1,3.45);

    % 水（低于水面的部分）
    \fill[gray!22] (0,0) -- (3,0) -- (4.2,0.9) -- (1.2,0.9) -- cycle;
    \fill[gray!22] (0,3) -- (3,3) -- (4.2,3.9) -- (1.2,3.9) -- cycle;

    % 棱柱
    \draw (A) -- (B) -- (C) -- (C1) -- (B1) -- (A1) -- cycle;
    \draw (B) -- (B1);
    \draw (C) -- (C1);
    \draw[dashed] (A1) -- (D1) -- (C1);
    \draw[dashed] (A) -- (D) -- (C);
    \draw[dashed] (D) -- (D1);

    % 内嵌正四棱锥（底面为 ABCD，顶点 P 在上）
    \draw[dashed] (A) -- (P);
    \draw[dashed] (B) -- (P);
    \draw[dashed] (C) -- (P);
    \draw[dashed] (D) -- (P);

    % 水面
    \draw[line width=0.8pt] (0,3) -- (3,3) -- (4.2,3.9) -- (1.2,3.9) -- cycle;

    \fill (P) circle (2pt);
    \node[right=1pt] at (P) {$P$};
    \node[below] at (2.1,-0.8) {图 1};
  \end{scope}

  % ---------- 图 2：倒置 ----------
  \begin{scope}[shift={(7.6,0)}]
    \coordinate (A2)  at (0,0);
    \coordinate (B2)  at (3,0);
    \coordinate (C2)  at (4.2,0.9);
    \coordinate (D2)  at (1.2,0.9);
    \coordinate (A3) at (0,5);
    \coordinate (B3) at (3,5);
    \coordinate (C3) at (4.2,5.9);
    \coordinate (D3) at (1.2,5.9);
    \coordinate (P2) at (2.1,1.55);

    % 水（低于水面的部分；倒置后棱锥顶点朝下）
    \fill[gray!22] (0,0) -- (3,0) -- (4.2,0.9) -- (1.2,0.9) -- cycle;
    \fill[gray!22] (0,1.1) -- (3,1.1) -- (4.2,2.0) -- (1.2,2.0) -- cycle;

    % 棱柱
    \draw (A2) -- (B2) -- (C2) -- (C3) -- (B3) -- (A3) -- cycle;
    \draw (B2) -- (B3);
    \draw (C2) -- (C3);
    \draw[dashed] (A3) -- (D3) -- (C3);
    \draw[dashed] (A2) -- (D2) -- (C2);
    \draw[dashed] (D2) -- (D3);

    % 内嵌正四棱锥（倒置后顶点 P2 朝下）
    \draw[dashed] (A3) -- (P2);
    \draw[dashed] (B3) -- (P2);
    \draw[dashed] (C3) -- (P2);
    \draw[dashed] (D3) -- (P2);

    % 水面
    \draw[line width=0.8pt] (0,1.1) -- (3,1.1) -- (4.2,2.0) -- (1.2,2.0) -- cycle;

    \fill (P2) circle (2pt);
    \node[right=1pt] at (P2) {$P$};
    \node[below] at (2.1,-0.8) {图 2};
  \end{scope}
\end{tikzpicture}
